% GENERATED FILE. Edit canonical lessons, then run npm run generate:books.
% canonical-source-hash: fnv1a64:5837cceb5723d2c6
% canonical-lessons: TE-C162-muddupettu, TE-C162-kaugilincu, TE-C162-atapattincu, TE-C162-pogadu, TE-C162-varnincu

\chapter{To kiss, To hug, To tease, To praise, To describe}
\label{ch:te-to-kiss-to-hug-to-tease-to-praise-to-describe}
\hlchaptermodality{162}

\begin{chapteropening}
\textbf{By the end of this chapter:} \emph{I can say to kiss, to hug, to tease, to praise and to describe.}

Complete the last lesson of chapter 162: I can say to kiss, to hug, to tease, to praise and to describe.
\end{chapteropening}

\section[\texorpdfstring{\te{ముద్దుపెట్టు} (muddupeṭṭu)}{ముద్దుపెట్టు (muddupeṭṭu)}]{\te{ముద్దుపెట్టు} (muddupeṭṭu) — to kiss}

\label{lesson:TE-C162-muddupettu}



\begin{quote}
Before the new one: say the Telugu for to feel shy, then the Telugu for to be surprised.
\end{quote}

\subsection*{You'll want to know: \te{ముద్దుపెట్టు}}
\textbf{\te{ముద్దుపెట్టు}} — \emph{muddupeṭṭu} — \textquotedblleft{}to kiss\textquotedblright{}.

\begin{grammarlens}[title={What you've built}]
One more word for this chapter.
\end{grammarlens}

\subsection*{Your turn}
\begin{itemize}
\raggedright
  \item \emph{Say it:} \emph{muddupeṭṭu}
  \item \emph{Say it:} \emph{muddupeṭṭu}, once more
  \item \emph{Say it:} say \emph{muddupeṭṭu}
  \item \emph{Recall:} say the Telugu for to feel shy, then the Telugu for to be surprised, then say \emph{muddupeṭṭu} again
\end{itemize}

\begin{tcolorbox}[breakable,colback=teal!4,colframe=teal!35!black,title={Before you move on}]
What does \te{ముద్దుపెట్టు} mean? (\textquotedblleft{}To kiss\textquotedblright{}.) Say it once more.
\end{tcolorbox}
\section[\texorpdfstring{\te{కౌగిలించు} (kaugiliñcu)}{కౌగిలించు (kaugiliñcu)}]{\te{కౌగిలించు} (kaugiliñcu) — to hug, to embrace}

\label{lesson:TE-C162-kaugilincu}



\begin{quote}
Before the new one: say the Telugu for to be surprised, then the Telugu for to kiss.
\end{quote}

\subsection*{You'll want to know: \te{కౌగిలించు}}
\textbf{\te{కౌగిలించు}} — \emph{kaugiliñcu} — \textquotedblleft{}to hug, to embrace\textquotedblright{}.

\begin{grammarlens}[title={What you've built}]
One more word for this chapter.
\end{grammarlens}

\subsection*{Your turn}
\begin{itemize}
\raggedright
  \item \emph{Say it:} \emph{kaugiliñcu}
  \item \emph{Say it:} \emph{kaugiliñcu}, once more
  \item \emph{Say it:} say \emph{kaugiliñcu}
  \item \emph{Recall:} say the Telugu for to be surprised, then the Telugu for to kiss, then say \emph{kaugiliñcu} again
\end{itemize}

\begin{tcolorbox}[breakable,colback=teal!4,colframe=teal!35!black,title={Before you move on}]
What does \te{కౌగిలించు} mean? (\textquotedblleft{}To hug, to embrace\textquotedblright{}.) Say it once more.
\end{tcolorbox}
\section[\texorpdfstring{\te{ఆటపట్టించు} (āṭapaṭṭiñcu)}{ఆటపట్టించు (āṭapaṭṭiñcu)}]{\te{ఆటపట్టించు} (āṭapaṭṭiñcu) — to tease}

\label{lesson:TE-C162-atapattincu}



\begin{quote}
Before the new one: say the Telugu for to kiss, then the Telugu for to hug.
\end{quote}

\subsection*{You'll want to know: \te{ఆటపట్టించు}}
\textbf{\te{ఆటపట్టించు}} — \emph{āṭapaṭṭiñcu} — \textquotedblleft{}to tease\textquotedblright{}.

\begin{grammarlens}[title={What you've built}]
One more word for this chapter.
\end{grammarlens}

\subsection*{Your turn}
\begin{itemize}
\raggedright
  \item \emph{Say it:} \emph{āṭapaṭṭiñcu}
  \item \emph{Say it:} \emph{āṭapaṭṭiñcu}, once more
  \item \emph{Say it:} say \emph{āṭapaṭṭiñcu}
  \item \emph{Recall:} say the Telugu for to kiss, then the Telugu for to hug, then say \emph{āṭapaṭṭiñcu} again
\end{itemize}

\begin{tcolorbox}[breakable,colback=teal!4,colframe=teal!35!black,title={Before you move on}]
What does \te{ఆటపట్టించు} mean? (\textquotedblleft{}To tease\textquotedblright{}.) Say it once more.
\end{tcolorbox}
\section[\texorpdfstring{\te{పొగడు} (pogaḍu)}{పొగడు (pogaḍu)}]{\te{పొగడు} (pogaḍu) — to praise}

\label{lesson:TE-C162-pogadu}



\begin{quote}
Before the new one: say the Telugu for to hug, then the Telugu for to tease.
\end{quote}

\subsection*{You'll want to know: \te{పొగడు}}
\textbf{\te{పొగడు}} — \emph{pogaḍu} — \textquotedblleft{}to praise\textquotedblright{}.

\begin{grammarlens}[title={What you've built}]
One more word for this chapter.
\end{grammarlens}

\subsection*{Your turn}
\begin{itemize}
\raggedright
  \item \emph{Say it:} \emph{pogaḍu}
  \item \emph{Say it:} \emph{pogaḍu}, once more
  \item \emph{Say it:} say \emph{pogaḍu}
  \item \emph{Recall:} say the Telugu for to hug, then the Telugu for to tease, then say \emph{pogaḍu} again
\end{itemize}

\begin{tcolorbox}[breakable,colback=teal!4,colframe=teal!35!black,title={Before you move on}]
What does \te{పొగడు} mean? (\textquotedblleft{}To praise\textquotedblright{}.) Say it once more.
\end{tcolorbox}
\section[\texorpdfstring{\te{వర్ణించు} (varṇiñcu)}{వర్ణించు (varṇiñcu)}]{\te{వర్ణించు} (varṇiñcu) — to describe}

\label{lesson:TE-C162-varnincu}



\begin{quote}
Before the new one: say the Telugu for to tease, then the Telugu for to praise.
\end{quote}

\subsection*{You'll want to know: \te{వర్ణించు}}
\textbf{\te{వర్ణించు}} — \emph{varṇiñcu} — \textquotedblleft{}to describe\textquotedblright{}.

\begin{grammarlens}[title={What you've built}]
One more word for this chapter.
\end{grammarlens}

\subsection*{Your turn}
\begin{itemize}
\raggedright
  \item \emph{Say it:} \emph{varṇiñcu}
  \item \emph{Say it:} \emph{varṇiñcu}, once more
  \item \emph{Say it:} say \emph{varṇiñcu}
  \item \emph{Recall:} say the Telugu for to tease, then the Telugu for to praise, then say \emph{varṇiñcu} again
\end{itemize}

\begin{tcolorbox}[breakable,colback=teal!4,colframe=teal!35!black,title={Before you move on}]
What does \te{వర్ణించు} mean? (\textquotedblleft{}To describe\textquotedblright{}.) Say it once more.
\end{tcolorbox}
